\scanpage{134}\mspage{68}

sur l'ens $S$ sera par définition la donnée d'une classe d'équivalence\inserted{ $\Pi$}, au sens de (18) de structures tr. sympl précongruentielles.

% slanted
\margin{$\Pi\hookrightarrow\mathfrak P(\mathfrak P_3(S))$ ou $\Pi\hookrightarrow\mathfrak P(\mathfrak P_2(S)$ \unsure{au|etc.} choix\ldots}

\unsure{A'|À} une telle structure est donc associée une suite exacte canonique (17), où $G_0$ est le groupe d'autom. commun des structures $(S,F_i)_{i\in\Pi}$, et où $Z$ est le groupe $Z_{F_i}$, commun pour tous les $i\in\Pi$, $\Pi$ étant de façon canonique un torseur sous $Z$.

\uline{NB} À vrai dire, cette axiomatique n'est pas encore réduite vraiment à l'os, en ce sens que les structures triangulées n'y jouent qu'un rôle logique assez accessoire -- les réflexions aboutissant aux axiomes To 3, To 4 (To 1 et To 2 $\unsure{y|{}}$ tombent) pouvant se développer le plus naturellement dans le context des \uline{graphes} plutôt que dans celui des structures triangulées qu'on peut associer à cer\uline{tains} graphe.

Pour en revenir quand $=$ aux structures triangulées, on voit que la structure \deleted{\unsure{p|P}}triangulée précongruentielle
\[
\Pi\in\mathfrak P(\mathfrak P(\mathfrak P_3(S)))\tag*{(19)}
\]

% archive sheet number: 133
